\documentclass[aps,onecolumn,nofootinbib,pra]{article}

\usepackage{../../spec_files/arxiv}
\usepackage{amsmath,amsfonts,amssymb,amsthm,enumerate,times}
\usepackage{mathtools}
\usepackage[usenames,dvipsnames]{color}
\usepackage{hyperref}
\hypersetup{
    breaklinks,
    colorlinks,
    linkcolor=gray,
    citecolor=gray,
    urlcolor=gray,
    pdftitle={},
    pdfauthor={Alex Goessmann}
}

\usepackage{tikz}
\usepackage{graphicx}
\usepackage{float}
\usepackage{comment}
\usepackage{csquotes}
\usepackage{braket}
\usepackage{listings}
\usepackage{verbatim}
\usepackage{etoolbox}
\usepackage[utf8]{inputenc}
\usepackage[english]{babel}
\usepackage[T1]{fontenc}
\usepackage{titlesec}
\usepackage{fancyhdr}
\usepackage{bbm}
\usepackage{bm}
\usepackage{algpseudocode}
\usepackage{algorithm}
\usepackage{lipsum}

\newtheorem{remark}{Remark}
\newtheorem{theorem}{Theorem}
\newtheorem{lemma}{Lemma}
\newtheorem{corollary}{Corollary}
\newtheorem{definition}{Definition}
\newtheorem{example}{Example}

\pretolerance=500
\tolerance=100
\emergencystretch=10pt

% Bibliography
\DeclareUnicodeCharacter{FB01}{fi}
\usepackage[round]{natbib}
\usepackage{wasysym}

\newcommand{\red}[1]{\textcolor{red}{#1}}

\usepackage{import}
\subimport{../../macros/}{all-macros}

\usepackage{minted}

\begin{document}
    \title{Tensorized Knowledge Units}

    \maketitle
    \date{\today}

    \begin{abstract}
        We investigate, whether tensor networks could serve as a knowledge representation format for the project Alexandria \cite{schuhmann_project_2025}.
        To represent uncertainties and relations between claims we apply the formalism discussed in \cite{goessmann_tensor_2026}.
    \end{abstract}


    \section{Incorporating uncertainties}

    \begin{center}
        \textit{
            Climate change is happening with $80\%$ probability.
        }
    \end{center}
    \stantikzpic{
        \draw (-2,-3) rectangle (2,1);
        \drawtextnode{0}{-1}{$\begin{pmatrix}
                                  0.2 \\ 0.8
        \end{pmatrix}$}
        \drawvwire{0}{-3}{$\catvariableof{0}$}
    }


    \section{Modelling relations between claims}

    \begin{center}
        \textit{
            When the polar ice is melting, then climate change is happening.
        }
    \end{center}

    \begin{itemize}
        \item $\catvariableof{0}$: Polar ice is melting
        \item $\catvariableof{1}$: Climate change is happening
    \end{itemize}



    \stantikzpic{
        \draw (-2,-3) rectangle (2,1);
        \drawtextnode{0}{-1}{$\begin{pmatrix}
                                  1 & 1 \\
                                  0 & 1
        \end{pmatrix}$}
        \drawvwire{-1}{-3}{$\catvariableof{0}$}
        \drawvwire{1}{-3}{$\catvariableof{1}$}
    }


    \section{Inference}


    \begin{center}
        \textit{
            Assuming that the polar ice is melting, what is the probability of climate change, according to the paper?
        }
    \end{center}

    This would be a contraction with $\tbasisat{\catvariableof{0}}$

    \stantikzpic{
        \draw (-2,-3) rectangle (2,1);
        \drawtextnode{0}{-1}{$\begin{pmatrix}
                                  1 & 1 \\
                                  0 & 1
        \end{pmatrix}$}
        \drawvwire{-1}{-3}{$\catvariableof{0}$}
        \drawvwire{1}{-3}{$\catvariableof{1}$}
        \drawsmalltensor{-1}{-6}{$\tbasis$}

        \drawtextnode{5}{-1}{$=$}

        \draw (8,-3) rectangle (12,1);
        \drawtextnode{10}{-1}{$\begin{pmatrix}
                                   0 \\ 1
        \end{pmatrix}$}

        \drawvwire{10}{-3}{$\catvariableof{1}$}
    }


    \section{Similarity by contractions}

    The extracted tensorized Knowledge Units can be contracted to observe the similarity between their claims.
    For example, two knowledge units with contradicting claim have vanishing contractions.

%
%    Further ideas:
%    \begin{itemize}
%        \item Usage of distance measures of probability distributions
%        \item Entailment of knowledge units
%    \end{itemize}



    \bibliographystyle{plainnat}
    \bibliography{../../references.bib}


\end{document}